\section{Evaluation}

We will answer the following research questions.  

\vspace*{3pt}
\noindent
{\bf RQ1: How effective is our approach in localizing concurrent bugs?}

\vspace*{3pt}
\noindent
{\bf RQ2: How does ConFL compare with state-of-the-art across different bug-report conditions?}

\vspace*{3pt}
\noindent
{\bf RQ3: What is the contribution of each component in our approach?}

\vspace*{3pt}
\noindent
{\bf RQ4:  How well does our approach generalize to unseen data and different LLMs?}

\subsection{Experimental Setup}


\subsubsection{Dataset}

The first dataset is derived from projects used in prior concurrent fault localization~\cite{shao2023information} and bug report classification studies~\cite{shao2026identifying}. As shown in Table~\ref{table:git1}, Dataset$_{Git}$ contains eight large-scale open-source systems, including druid, grpc-java, presto, pulsar, redisson, rocketmq, trino, and vert.x, with 4--100 concurrency bug reports per project. The projects range from roughly 9K to over 700K methods and cover diverse concurrency styles, including lock-based synchronization, executor-based concurrency, asynchronous callbacks, and mixed interaction patterns. Ground-truth faulty methods are identified from bug-fixing commits.



To evaluate generalization beyond potential training-data exposure, we additionally construct a post-cutoff dataset containing concurrency bugs reported after the LLM training cutoff date. We follow the same collection procedure as Dataset$_{Git}$ and retain only bugs with concurrency-related descriptions, associated fix commits, and verifiable faulty methods. Since several Dataset$_{Git}$ projects yielded few eligible bugs, we further include widely used Java systems frequently used in concurrency-related research, including Camel, Cassandra, Hadoop, and HBase~\cite{shao2023information, shao2026identifying}. As shown in Table~\ref{table:git2}, the resulting dataset contains 32 bugs in total.



\begin{table*}[t]\footnotesize
\centering
\begin{minipage}{0.48\linewidth}
\centering
\caption{Dataset$_{Git}$}
\label{table:git1}
\begin{tabular}{c|ccc}
\toprule
Project & \#BR & Avg.\#M & Avg.\#F \\ \hline
druid  & 31      & 22,692      &  2,954        \\ 
grpc-java    &  43     & 9,970      & 643       \\ 
presto    &  7     &  23,350     &  2,790         \\ 
pulsar    &  72     &  27,917     & 2,998          \\ 
redisson    & 100       &  14,262     & 1,216          \\ 
rocketmq    & 4      & 15,123      & 1,589         \\ 
trino    & 17      &  708,356     & 8,160          \\ 
vert.x     & 48      &  11,525     & 761    \\ 
\textbf{total}     & 322      &  14,251  &  1,808   \\ 
\bottomrule
\end{tabular}
\end{minipage}
\hfill
\begin{minipage}{0.48\linewidth}
\centering
\caption{Dataset$_{Post}$}
\label{table:git2}
\begin{tabular}{c|ccc}
\toprule
Project & \#BR & Avg.\#M & Avg.\#F \\ \hline
druid  & 1      & 73,782      & 8,608         \\ 
grpc-java    & 4      & 20,839      & 1,512       \\ 
pulsar    & 13      & 37,971      &  4,200         \\ 
redisson    & 3      & 31,777      & 2,609          \\ 
camel    & 4      & 180,000      & 22,406          \\ 
cassandra    & 2      & 42,408     & 3,463          \\ 
hadoop     & 1      & 136,170      & 11,829    \\ 
hbase     & 4      &  69,600     & 5,035    \\ 
\textbf{total}    & 32      & 72,945      &  6,425   \\ 
\bottomrule
\end{tabular}
\end{minipage}
\end{table*}
\vspace{-0.5em}


\subsubsection{Baseline Selection.}

Our framework operates strictly in a bug-report–only setting. Therefore, we select baselines that can be executed under the same constraint and that represent the strongest available techniques in their respective categories.

\noindent
\textbf{IR-Based Baselines.}
Most IR-based FL techniques work at the file level and are not directly applicable to method-level localization. We include two representative methods: \textbf{BRTracer} \cite{BTrace} — shown in prior studies \cite{shao2024enhancing, lee2018bench4bl} to be one of the strongest IRFL baselines. We adapt it to method-level by treating each method as an independent document. \textbf{BoostNSift} \cite{razzaq2021boostnsift} — a recent IR-based method designed specifically for method-level FL, providing a competitive IR baseline for our evaluation. These baselines represent the upper bound of lexical-similarity approaches under bug-report–only conditions. We additionally include \textbf{BLCoiR}~\cite{shao2023information}, the only prior bug-report–based fault localization technique specifically designed for concurrent bugs. Since BLCoiR originally operates at file-level granularity, we similarly adapt it to method-level retrieval by treating each method as an independent document. For file-level comparison, we aggregate \Name{}'s method-level outputs using a max-score rule, where the score of a file is defined as the maximum score among its methods. Together, these baselines represent strong lexical-similarity approaches under bug-report–only settings.


\noindent
\textbf{LLM-Based Baselines.}
Many recent LLM-based approaches (e.g., AutoFL, AgentFL) rely on test-case execution and cannot run without dynamic information. As discussed in our motivation, concurrent bugs rarely have reliable test cases, making these methods incompatible with our evaluation setting.
We therefore select: \textbf{FlexFL} — a recent LLM-based FL framework that supports bug-report–only mode and thus can be fairly compared with ours.

\noindent
\textbf{Other Baselines.}
Supervised FL techniques such as DreamLoc~\cite{dreamloc} and rLocator~\cite{locator} require thousands of labeled bugs per project, while our projects contain only 4–40 real-world concurrent bugs, making such training infeasible.
Other learning-based methods like DeepFL depend on mutation-based synthetic bugs and large test suites, which are unavailable for concurrent bug scenarios. Thus, these approaches cannot serve as practical baselines.

Recent LLM  methods for Python (e.g., CoSIL) rely on building full program graphs and performing LLM-guided exploration. These techniques do not scale to Java due to: call graphs with hundreds of thousands to millions of nodes,	extensive virtual dispatch, and lack of existing Java adaptations.

Spectrum-based and hybrid techniques (e.g., Tarantula~\cite{Tarantula}, Ochiai~\cite{abreu2009practical}, DStar~\cite{dst}, DeepFL~\cite{dl1}, Grace~\cite{dl2}) rely on failing/passing tests and coverage information, which are difficult to obtain reliably for concurrency bugs due to nondeterministic scheduling. A concurrency test may pass or fail across runs without code changes, weakening coverage-based statistical correlation. Supervised learning-based methods also require large amounts of labeled training data, which is impractical given the limited number of concurrency bugs available per project in our benchmark (4--100 bugs/project).

\subsubsection{Performance Metrics.}

We evaluate \Name{} using three standard metrics widely adopted in prior fault localization studies~\cite{flexfl, autofl,yu2008empirical,jones2005empirical,chilimbi2009holmes,santelices2009lightweight}.

\textbf{Top@K} measures the proportion of bug reports whose faulty method appears within the top-$k$ results ($k\!\in\!\{1,5,10\}$), defined as
$Top@K = |R_k| / n$, where $|R_k|$ is the number of reports with at least one correct faulty method ranked in the top-$k$, and $n$ is the total number of bug reports.

\textbf{MRR} (Mean Reciprocal Rank) captures how early the first faulty method appears:
$MRR = \frac{1}{n}\sum_{j=1}^{n} \frac{1}{rank_j}$,
where $rank_j$ is the rank of the first relevant faulty method for the $j$-th report.

\textbf{MAP} (Mean Average Precision) evaluates the overall ranking quality of all faulty methods:
$MAP = \frac{1}{n}\sum_{j=1}^{n} AvgP_j$,
where $AvgP_j = \frac{1}{|K_j|}\sum_{k\in K_j} Prec@k$ and
$Prec@k = \frac{1}{k}\sum_{i=1}^{k} IsRelevant(i)$.
Here, $K_j$ denotes the set of true faulty methods for report $j$, and $IsRelevant(i)=1$ if the $i$-th ranked method is faulty. Higher values indicate better localization performance.


\subsection{Results Analysis}

\subsubsection{RQ1: How effective is our approach in localizing concurrent bugs?}

Table~\ref{table:overall} presents the fault localization performance of our approach on DatasetGit.
ConFL achieves 132 Top-1 hits, 186 Top-3 hits, and 216 Top-5 hits, substantially outperforming all baselines.
In contrast, the strongest baseline, FlexFL, achieves 82 Top-1 hits, 110 Top-3 hits, and 123 Top-5 hits, while IR-based techniques perform significantly worse.

In terms of ranking quality, ConFL attains an MRR of 0.503 and a MAP of 0.486, improving upon FlexFL by +0.189 MRR and +0.178 MAP.
Compared to traditional IR-based approaches (BRTracer and BoostNSift), the improvements are even more pronounced, exceeding +0.40 in both MRR and MAP.
These results indicate that ConFL not only identifies the correct faulty methods more frequently, but also ranks them more consistently near the top of the candidate list.

\vspace{-8pt}
\begin{table}[H]\footnotesize
\caption{Performance on Dataset$_{Git}$} \label{table:overall}
\vspace{-8pt}
\begin{tabular}{c|ccc|cc}
\toprule
Technique & Top-1 & Top-3 & Top-5 & MRR & MAP \\ \hline
BRTracer  &  18     &  32     &  47     & 0.084    & 0.082     \\ 
BoostN    &   24    &  41     &   50    & 0.103    &  0.099   \\ 
BLCoiR  &  5     &  11    &  20     & 0.030    & 0.030     \\ \hline
Flexfl    &  82     &  110     &  123     &  0.314   & 0.308    \\ 
ConFL     &  132     &  186     &  216     & \textbf{0.503}    & \textbf{0.486}    \\ \bottomrule
\end{tabular}
\end{table}
\vspace{-8pt}



\noindent
\textbf{Comparison with BLCoiR.} BLCoiR performs poorly at the method level, achieving only 0.030 MRR and 0.030 MAP. This is mainly due to a granularity mismatch between its sparse, file-level query design and method-level retrieval. BLCoiR retains only coarse-grained concurrency keywords and class-level information, which are effective for file-level retrieval where files contain rich aggregated context, but insufficient for individual methods that are themselves sparse and localized. To provide a fairer comparison, we additionally evaluate both techniques at file-level granularity, following the original BLCoiR setting. As shown in Table~\ref{tab:file}, \Name{} still substantially outperforms BLCoiR, improving MRR from 0.32 to 0.76 and MAP from 0.21 to 0.70. These results suggest that \Name{}'s gains primarily come from structured concurrency interaction modeling rather than localization granularity alone.

\begin{wraptable}{r}{0.3\textwidth} \footnotesize
\vspace{-10pt}
\centering 
\caption{File-level comparison.}
\vspace{-8pt}
\label{tab:file}
\begin{tabular}{lcc}
\toprule
Technique & MRR & MAP \\
\midrule
BLCoiR & 0.32 & 0.21 \\
ConFL & \textbf{0.76} & \textbf{0.70} \\
\bottomrule
\end{tabular}
\vspace{-10pt}
\end{wraptable}

Overall, the strong gains in Top-$k$ accuracy, MRR, and MAP on Dataset$_{Git}$ demonstrate the effectiveness of ConFL in localizing real-world concurrent bugs using bug reports alone, even in large and diverse codebases mined from GitHub.


\noindent\textbf{Evaluating Explanation Quality.}
As described in Section~4.3, ConFL explanations identify suspicious methods, thread entry points, shared resources, and synchronization context. The goal of this evaluation is to assess whether these explanations faithfully describe the concurrency behavior of the localized fault rather than their fluency or writing quality.

\noindent\textbf{Evaluation Setup.}
Automated metrics such as BLEU or BERTScore cannot determine whether claims about thread interactions are factually correct, and no benchmark of ground-truth concurrency explanations exists. We therefore conduct manual evaluation against the source code and developer fix commits. Annotators classify each explanation as \textbf{Faithful} (all concurrency claims are supported by the code), \textbf{Partial} (the explanation is generally correct but omits relevant interactions), or \textbf{Hallucinated} (the explanation contains unsupported thread interactions or synchronization behavior).

We evaluate only successfully localized bugs, since explanations based on incorrect localizations are not meaningful. We randomly sample 50 bugs and generate three explanations per bug under two settings: \textit{w/ DSL}, where the LLM receives interaction DSLs, and \textit{w/o DSL}, where the LLM receives only raw source code. This yields 300 explanations in total. Two annotators independently labeled all explanations, achieving 88.3\% agreement with Cohen's $\kappa = 0.767$, indicating substantial agreement.

\begin{wraptable}{r}{0.32\textwidth}\footnotesize
\vspace{-10pt}
\centering
\caption{Explanation quality.}
\vspace{-8pt}
\label{tab:explanation}
\begin{tabular}{lcc}
\toprule
Category & w/ DSL & w/o DSL \\
\midrule
Faithful     & \textbf{78\%} & 70\%          \\
Partial      & 10\%          & 5\%           \\
Hallucinated & \textbf{12\%} & 25\%          \\
\bottomrule
\end{tabular}
\vspace{-10pt}
\end{wraptable}


\noindent\textbf{Results.}
Table~\ref{tab:explanation} shows that the interaction DSL improves the faithful rate from 70\% to 78\% while reducing hallucinations from 25\% to 12\%. This reduction is particularly important because hallucinated thread interactions or synchronization behavior can mislead developers during debugging. The improvement stems from the DSL grounding explanations in explicit thread entry points, shared variables, and access relationships extracted from static analysis, thereby constraining the LLM to reason over verifiable concurrency structure.

\noindent\textbf{Illustrative Case.}
The motivating example \textbf{Druid-3174} illustrates this effect. Without the DSL, the LLM explains only local behavior inside \texttt{RemoteTaskRunner}. With the DSL, ConFL additionally captures that \texttt{WorkerTaskMonitor} initiates concurrent execution while \texttt{RemoteTaskRunner} launches additional threads accessing shared task-state resources, enabling the explanation to describe the fault in terms of concrete cross-thread interactions.

\vspace{-0.5em}
\vspace*{3pt}
\begin{tcolorbox}
\stepcounter{findingsCounter}
{\bf RQ\thefindingsCounter{} summary:} {
The results show that ConFL significantly outperforms all baselines on Dataset$_{Git}$, achieving higher Top-$k$ accuracy as well as substantially better MRR and MAP.
These gains indicate that ConFL not only localizes faulty methods more effectively, but also ranks them more reliably near the top.
Moreover, by providing interaction-level concurrency context, ConFL enables more informative explanations than method-level outputs alone, supporting both accurate and explainable fault localization using bug reports only.
}
\end{tcolorbox}


\subsubsection{RQ2: How does ConFL compare with state-of-the-art across different bug-report conditions?}

While RQ1 establishes that ConFL outperforms both IR-based baselines (BRTracer, BoostNSift) and the LLM-based baseline FlexFL on aggregate metrics, it does not reveal where this improvement comes from or under what conditions each technique succeeds or fails. To identify the source of ConFL's gains, we stratify the evaluation by bug-report quality, comparing all techniques under three reporting conditions that reflect different levels of structural guidance, following prior fault-localization studies~\cite{shao2024enhancing, blizzard, shao2023information}.
Specifically, we consider: (i) PE reports that include explicit identifiers aligned with the ground truth (PE with hint),
(ii) PE reports that contain misleading or irrelevant identifiers (PE without hint), and
(iii) NL-only reports that provide only natural-language descriptions without code identifiers.


Table~\ref{tab:db} summarizes the results.
When strong structural hints are present (PE with hint), all LLM-based approaches benefit substantially.
In this setting, ConFL achieves the best performance (MRR = 0.756, MAP = 0.724), outperforming FlexFL by +0.15 MRR and +0.13 MAP, and significantly surpassing IR-based baselines.
This result indicates that even with accurate identifiers, concurrency-aware static context provides complementary signals beyond surface-level matching.

The \emph{PE without hint} setting is the most challenging, as misleading program-entity mentions can actively steer LLM reasoning toward incorrect code regions.
Under this condition, FlexFL suffers a severe performance collapse (MRR = 0.021, MAP = 0.021), performing comparably to or worse than IR-based methods.
In contrast, ConFL degrades more gracefully, retaining substantially higher performance (MRR = 0.189, MAP = 0.189).
Although noisy identifiers may still introduce false positives, hierarchy-based retrieval prevents the LLM from overfitting to individual symbols and enables more stable localization by grounding reasoning in concurrency-relevant structure.

Interestingly, NL-only reports are often easier to handle than \emph{PE without hint} reports.
While NL-only reports provide less explicit information, they do not actively mislead the reasoning process.
As a result, ConFL achieves higher performance on NL-only reports (MRR = 0.320, MAP = 0.321) than on PE without hint reports (MRR = 0.189, MAP = 0.189), as hierarchy-based retrieval can explore concurrency-relevant components and thread entries in a more unbiased manner.

\vspace{-8pt}

\begin{table*}[h]
\centering \footnotesize
\caption{MRR and MAP results under different bug report qualities.}\label{tab:db}
\vspace{-8pt}
\begin{subtable}{0.32\textwidth}
\centering
\caption{PE (with hint)}
\begin{tabular}{lcc}
\toprule
Technique & MRR & MAP \\
\midrule
BRTracer & 0.136 & 0.134 \\
boostNSift & 0.167 & 0.161 \\
BLCoiR & 0.048 & 0.048 \\
FlexFL & 0.607 & 0.597 \\
\Name{} & \textbf{0.756} & \textbf{0.724} \\
\bottomrule
\end{tabular}
\end{subtable}
\hfill
\begin{subtable}{0.32\textwidth}
\centering
\caption{PE (without hint)}
\begin{tabular}{lcc}
\toprule
Technique & MRR & MAP \\
\midrule
BRTracer & 0.026 & 0.024 \\
boostNSift & 0.042 & 0.040 \\
BLCoiR & 0.007 & 0.007 \\
FlexFL & 0.021 & 0.021 \\
\Name{} & \textbf{0.189} & \textbf{0.189} \\
\bottomrule
\end{tabular}
\end{subtable}
\hfill
\begin{subtable}{0.32\textwidth}
\centering
\caption{NL}
\begin{tabular}{lcc}
\toprule
Technique & MRR & MAP \\
\midrule
BRTracer & 0.052 & 0.054 \\
boostNSift & 0.054 & 0.049 \\
BLCoiR & 0.018 & 0.020 \\
FlexFL & 0.023 & 0.021 \\
\Name{} & \textbf{0.320} & \textbf{0.321} \\
\bottomrule
\end{tabular}
\end{subtable}
\end{table*}


\vspace{-8pt}

Overall, these results show that although LLM-based techniques can perform well when guided by accurate program-entity hints, they are highly sensitive to noise and misleading information.
By explicitly modeling concurrency interactions and decoupling localization from unreliable lexical signals, ConFL achieves both superior accuracy and improved robustness under realistic and noisy bug-report conditions.

\noindent
\textbf{Case Study.} We revisit the motivating example druid-1360, to briefly illustrate how our framework performs hierarchical fault localization and produces explainable results in practice.
Starting from the bug report, the framework first identifies the \texttt{indexing-service} component and further narrows the search space to the package \texttt{io.druid.indexing.overlord.autoscaling}.
Within this scope, \texttt{SimpleResourceManagement} is identified as a key thread entry point involved in autoscaling decisions, which interacts with other concurrent execution contexts such as the \texttt{RemoteTaskRunner}.

By constructing interaction-level DSL representations for these thread entries, the framework exposes concurrent accesses to shared task and worker state.
Based on this interaction-level context, the LLM identifies a race between task provisioning and termination, localizing the fault to \texttt{doProvision} and \texttt{doTerminate} in \texttt{SimpleResourceManagement}.
This example demonstrates how hierarchical retrieval and interaction-level modeling guide the LLM to focus on a small set of concurrency-critical methods, enabling concise and concurrency-aware explanations without exploring the entire codebase.

\vspace{-0.5em}
\vspace*{3pt}
\begin{tcolorbox}
\stepcounter{findingsCounter}
{\bf RQ\thefindingsCounter{} summary:} {
Across all bug-report conditions, ConFL consistently outperforms both IR-based and LLM-based baselines.
While LLM-based methods benefit from accurate program-entity hints, they are highly sensitive to noisy or misleading identifiers.
By grounding reasoning in structured concurrency knowledge and hierarchy-based retrieval, ConFL achieves superior robustness and maintains strong localization performance even under misleading or identifier-free bug reports.
These results demonstrate that explicitly modeling concurrency interactions is critical for reliable fault localization in realistic settings.
}
\end{tcolorbox}


\subsubsection{RQ3: What is the contribution of each component in our framework?}

To quantify the contribution of each major component in ConFL, we conduct an ablation study by disabling one component at a time.
Specifically, we evaluate three variants:
(i) \textbf{w/o Hierarchy}, which removes hierarchical retrieval and directly performs flat retrieval over the entire search space;
(ii) \textbf{w/o CKB}, which disables the concurrency knowledge base and provides only basic method-level information without structured concurrency semantics; and
(iii) \textbf{w/o DSL}, which removes the interaction-level DSL representation and feeds unstructured concurrency information to the LLM.
The full system (\textbf{ConFL}) serves as the reference.

\vspace{-0.5em}
\vspace{-4pt}
\begin{table}[H]\footnotesize
\caption{Ablation study} \label{table:ablation}
\vspace{-8pt}
\begin{tabular}{c|ccc|cc}
\hline \toprule
Variant & Top-1 & Top-3 & Top-5  & MRR & MAP \\ \hline
w/o Hierarchy  &  92     & 140      &  158 & 0.356   & 0.345     \\ 
w/o CKB    & 104      & 157      & 178     & 0.405   & 0.391   \\ 
w/o DSL    & 117      & 170      & 193     & 0.449   & 0.432     \\ \hline
ConFL     & 132      & 186      & 216    & \textbf{0.503}    & \textbf{0.486}     \\ \bottomrule
\end{tabular}
\end{table}
\vspace{-0.5em}




\begin{wrapfigure}{r}{0.35\textwidth} 
    \centering 
    \includegraphics[width=0.35\textwidth]{figs/h.pdf}
    \caption{Accuracy of Hierarchy}\label{fig:hierarchy-accuracy}
\end{wrapfigure}

As shown in Table~\ref{table:ablation}, removing any component leads to a clear degradation in localization performance, confirming that all components contribute meaningfully to the effectiveness of ConFL.
Among the ablated variants, \textbf{w/o Hierarchy} suffers the largest performance drop, with Top-1 accuracy decreasing from 132 to 92.
This result highlights the critical role of hierarchical retrieval in progressively constraining the search space and guiding the LLM toward relevant concurrency contexts.
Without hierarchy, the model is more likely to explore irrelevant code regions and follow incorrect reasoning paths.

The \textbf{w/o CKB} variant also shows a substantial decline (Top-1 = 104), indicating that the concurrency knowledge base provides essential structural and behavioral signals beyond method signatures.
Without explicit knowledge of shared variables, locks, and thread-entry properties, the LLM’s ability to reason about concurrency interactions is significantly weakened.



In contrast, the \textbf{w/o DSL} variant exhibits a smaller but still noticeable performance drop (Top-1 = 117).
This suggests that while raw concurrency information is useful, the structured interaction-level DSL plays an important role in organizing this information into a form that is easier for the LLM to reason over.
By explicitly encoding thread entries, shared resources, and access patterns, the DSL improves ranking accuracy and overall localization quality.

Overall, the full ConFL framework consistently achieves the best Top-k results, demonstrating that hierarchical retrieval, the concurrency knowledge base, and the interaction-level DSL are complementary and jointly necessary for effective and explainable concurrent fault localization.

% \sh{R2Q2 search space reduction is not the reason}

% \shuai{The CKB is designed to preserve concurrency-relevant contexts rather than aggressively prune the program space. In practice, it still retains approximately 20\% of all methods, corresponding to 1,000--3,000 candidate methods per project. Moreover, existing fault localization approaches already implicitly narrow the search space using bug-report signals such as mentioned classes or stack traces. ConFL improves localization performance not merely through search space reduction, but by explicitly modeling structured cross-thread interactions and enabling LLM-based reasoning over these interaction contexts.}


\paragraph{Accuracy of Hierarchy Retrieval.}
Figure~\ref{fig:hierarchy-accuracy} reports the accuracy of hierarchical retrieval at each stage of our framework.
Component-level accuracy measures whether the correct component is identified from the bug report, followed by package-, thread-entry-, and method-level accuracy measured conditionally at each subsequent stage.

\begin{wraptable}{r}{0.3\textwidth} \footnotesize
\vspace{-10pt}
\centering 
\caption{BoostN with CKB} \label{tab:boost}
\begin{tabular}{lcc}
\toprule
Technique & MRR & MAP \\
\midrule
BoostN &  0.103  & 0.099 \\
BoostN (CKB) &  0.167   &  0.165 \\
\bottomrule
\end{tabular}
\vspace{-10pt}
\end{wraptable}

As shown in the figure, ConFL achieves consistently high accuracy across all hierarchy levels, exceeding 88\% at every stage.
In particular, component-level and thread-entry retrieval exceed 91\% accuracy, indicating that both structural localization and concurrency-entry identification are highly reliable.
Accuracy decreases slightly at deeper levels, such as package and method selection, which is expected as the search space becomes more fine-grained.



This strong performance reflects the high-cohesion and low-coupling properties of software systems.
Components and packages typically encapsulate related functionality and concurrency responsibilities, enabling hierarchy-guided retrieval to progressively narrow the search space while preserving correct localization decisions.
Overall, these results demonstrate that hierarchical retrieval effectively leverages program modularity to support reliable interaction-level reasoning and LLM-based fault localization.


We emphasize that the gains of ConFL do not come solely from search-space reduction. Restricting retrieval to the CKB improves IR-based localization: as shown in Table~\ref{tab:boost}, BoostN(CKB) improves over BoostN from 0.103 to 0.167 MRR and from 0.099 to 0.165 MAP. However, the performance remains substantially lower than ConFL, indicating that filtering irrelevant methods alone is insufficient for concurrent fault localization. Similarly, the ablation study (Table~\ref{table:ablation}) shows that removing hierarchy-guided reasoning reduces Top-1 from 132 to 92 even under the same CKB setting. These results suggest that ConFL's gains mainly come from structured concurrency interaction modeling and hierarchy-guided reasoning rather than search-space reduction alone.



\vspace{-0.5em}
\vspace*{3pt}
\begin{tcolorbox}
\stepcounter{findingsCounter}
{\bf RQ\thefindingsCounter{} summary:} {
The ablation study confirms that hierarchical retrieval, the concurrency knowledge base, and the interaction-level DSL are all essential components of ConFL.
Hierarchical retrieval has the largest impact, as it effectively constrains the search space and guides the LLM toward relevant concurrency contexts.
This is further supported by the hierarchy-accuracy analysis, which shows consistently high accuracy across all levels, reflecting the high-cohesion and low-coupling structure of real-world software systems.
}
\end{tcolorbox}
\subsubsection{RQ4: How well does our approach generalize to unseen data and different LLMs?}

To assess the generalization ability of our framework, we evaluate two complementary dimensions:Performance on unseen (post cut-off) bugs and Performance across different LLM backbones. These experiments test whether our system overfits to training-era bug reports, and whether the gains from our concurrency-aware design persist across model architectures.

\vspace{-0.5em}
\begin{table}[H]\footnotesize
\caption{Performance on Dataset$_{Post}$} \label{tab:postcutoff}
\vspace{-8pt}
\begin{tabular}{c|ccc|cc}
\toprule
Technique & Top-1 & Top-3 & Top-5 & MRR & MAP \\ \hline
Flexfl    & 14      &   18    & 18       & 0.490   & 0.477    \\ \hline
ConFL     & 20      &  25     & 26      & \textbf{0.615}    & \textbf{0.598}    \\ \bottomrule
\end{tabular}
\end{table}
\vspace{-0.5em}


\noindent
\textbf{Generalization to Unseen Bugs.}
We evaluate our method on a post–cutoff dataset consisting of real-world concurrent bugs introduced after the LLM training cutoff, ensuring that none of the underlying models have prior exposure to these cases. As shown in Table~\ref{tab:postcutoff}, our framework maintains strong performance, achieving MRR = 0.615 and MAP = 0.598, substantially outperforming the LLM-based baseline FlexFL (MRR = 0.490, MAP = 0.477) as well as all IR-based methods, which fail to generalize in this setting. Notably, ConFL achieves 20 Top-1 hits out of 31, compared to FlexFL’s 14, demonstrating that our hierarchical retrieval and concurrency-aware reasoning remain effective even when applied to previously unseen systems and bug patterns. These results confirm that our approach generalizes well beyond the training-era distribution and remains robust under realistic post–cutoff conditions.

\noindent
\textbf{Generalization across different LLMs.}
We further evaluate the robustness of our framework using multiple LLM backbones, including GPT-4o, GPT-4o-mini, GPT-5.2 from OpenAI, and Claude~4.5 from Anthropic. Figure~\ref{fig:llm-generalization} shows that ConFL achieves its best performance with GPT-4o, while GPT-4o-mini delivers comparable results with only modest degradation. GPT-5.2 and Claude~4.5 exhibit slightly lower performance in some cases. Across all evaluated models, the MRR ranges from 0.41 to 0.503, corresponding to an absolute variation of 0.093, indicating relatively stable behavior across different LLM backbones.

\begin{wrapfigure}{r}{0.35\textwidth}
    \centering
    \includegraphics[width=0.35\textwidth]{figs/llms.pdf}
    \caption{Performance across LLMs} \label{fig:llm-generalization}
    \vspace{-8pt}
\end{wrapfigure}


Interestingly, newer or larger models do not consistently improve localization performance. We hypothesize that this behavior is related to differences in reasoning style and adherence to the hierarchy-guided retrieval process. Models such as GPT-4o and GPT-4o-mini appear to follow the structured interaction-level context more consistently, leading to more stable localization results.

Overall, these results suggest that ConFL benefits primarily from structured concurrency interaction modeling and hierarchy-guided reasoning rather than relying on a specific LLM architecture or unconstrained reasoning capability. The consistent improvements across different LLM backbones further demonstrate that the effectiveness of our framework stems from its structured reasoning process, making it robust and broadly applicable to a variety of foundation models.

\vspace{-0.5em}
\vspace*{3pt}
\begin{tcolorbox}
\stepcounter{findingsCounter}
{\bf RQ\thefindingsCounter{} summary:} {
The results show that ConFL generalizes well to both unseen concurrent bugs and different LLM backbones.
It consistently outperforms state-of-the-art baselines on post–cutoff bugs, indicating that its effectiveness does not rely on memorization.
Across LLMs, ConFL remains robust, with GPT-4o and GPT-4o-mini achieving the most stable performance.
Overall, these findings suggest that ConFL’s gains mainly come from hierarchy-guided retrieval and concurrency-aware interaction modeling, rather than from any specific LLM architecture.
}
\end{tcolorbox}
